% Section 6 — Limitations
% Status: v2, compressed to ~1.0–1.5 LNCS pages.
% Retained items:
%   (1) Revisit-heavy long-narration ceiling
%   (2) One honest spatial-axis failure (BurlingameDT4)
%   (3) Host-CPU latency scope
%   (4) No spatial-only baseline
%   (5) Multi-corpus coverage scope
%   (6) HDD is a systems/self-supervised result (no QA labels); F-HDD-3
%       geography-vs-visual confound (k-NN-cell control); ~3s GPS skew
%       corrected via analysis-time realign, source fix not yet re-extracted

\section{Limitations}
\label{sec:limitations}

\paragraph{Benchmark ceiling and coverage.}
On the Nymeria engine-check session, both eviction-free \psm{} and brute-force
CLIP reach only 13.4\% Hit@5: 162/187 action narrations occur in repeatedly
entered rooms with near-identical frames. A single annotated interval also
penalizes valid revisits. More importantly, the 14-session allocation study uses
third-person action narrations for Nymeria and target-frame-generated questions
for SLOPER4D/LookOut, not human-authored look-back requests. It therefore tests
retention under controlled retrieval proxies, not end-to-end wearable QA.

\paragraph{One honest spatial-axis failure.}
Of the 14 street-scale sessions, \texttt{BurlingameDT4} is the only one
where all three encoders fail the resolution-lift criterion: a visually
homogeneous parking-lot trajectory whose queries collapse onto a few
indistinguishable cells.

\paragraph{Evaluation scope.}
\psm's index latency and RSS are reported on host at the target CLIP-L
configuration ($d{=}768$, $R{=}128$, $C{=}60$; \S\ref{sec:results-memory-latency});
a prior Galaxy S22 figure was measured at a reduced ($d{=}128$, $R{=}4$) config and
is withdrawn as a deployment number, so on-device figures at the true configuration
remain to be measured. The CLIP encoder is expected to dominate on-device cost;
a distilled encoder (MobileCLIP~\cite{vasu2024mobileclip} or an AM-RADIO/E-RADIO
agglomerative model~\cite{ranzinger2024radio}) is the natural drop-in since \psm{}
is encoder-agnostic (\S\ref{sec:results-multi-corpus}). Our fixed-budget study
\emph{does} include coordinate-permutation and grid translation/rotation controls
(\S\ref{sec:results-main}), which isolate place-specific structure from generic
partition balancing; a $k$-nearest-cell / VPR-descriptor control remains future
work.

\paragraph{HDD is a memory-model characterization, not an engine-backed result.}
HDD ships no look-back QA labels, and the corpus-scale CLIP features needed to run
cardinality/cross-session retrieval \emph{through the C engine} are not available
(only a sanity drive was extracted). We therefore use HDD only to characterize
\psm{}'s per-cell memory model from the RTK trajectory (\S\ref{sec:results-hdd}),
as modeled logical state, and \emph{withdraw} the earlier corpus-scale cardinality
and cross-session-AUC numbers (computed by offline Python substitutes, not the
engine). Engine-backed multi-session HDD cardinality/retrieval --- with a
$k$-nearest-cell hard-negative control and purpose-built VPR
descriptors~\cite{keetha2023anyloc,alibey2023mixvpr} --- is deferred until the
drive-level features are extracted.
